% Gesamttest «Hydrostatik» — Quelle des PDFs (python3 scripts/build-lp-pdf.py hydrostatik).
% Musterlösung und Raster: bewertungspaket.tex. Alle Werte mit python3 nachgerechnet (05.10.2026).
% Jede Aufgabe fragt anders als Aufgaben, Übungen, Kontrollfragen und Leisten — G1 aus einem Diagramm
% (ablesen und rückwärts), G2 rückwärts (Tiefe aus dem Tauchcomputer im Bergsee), G3 Pascals Fassversuch
% (hydrostatisches Paradoxon), G4 Wagenheber rückwärts (Kolbenfläche aus der Handkraft), G5 Auftrieb in
% Luft (Ballon), G6 Aräometer (Dichte aus der Eintauchtiefe) — HOWTO §9.
\documentclass[11pt]{article}
\usepackage{../lp-druck}
\renewcommand{\lpname}{Hydrostatik}
\renewcommand{\lpteilgebiet}{4.5}
\renewcommand{\lpfuss}{Gesamttest · 25 Punkte}
% Freies Feld für Skizzen (ohne Linien)
\newcommand{\skizzenfeld}[1]{\par\vspace{4pt}\noindent\begin{tikzpicture}\draw[lplinie,dashed] (0,0) rectangle (\linewidth,#1);\end{tikzpicture}\par}
\begin{document}
\lpkopf{Gesamttest}{%
  \vspace{4pt}\noindent\begin{tabularx}{\linewidth}{@{}X X@{}}
  Code (kein Name): \hrulefill & Punkte: \hrulefill\ / 25
  \end{tabularx}}

\begin{lpinfo}{Anleitung}
\textbf{Zeit:} rund 30 Minuten. \textbf{Hilfsmittel:} Taschenrechner und Formelsammlung, in allen Teilen.
Schreib zu jeder Aufgabe zuerst die Formel, dann die Werte \textbf{mit Einheiten}. Punkte gibt es für den Weg,
für das Ergebnis mit Einheit, für Skizzen und für Begründungen in eigenen Worten. Rechne mit $g = 9.81\;\text{m/s}^2$ und
$\rho_\text{Wasser} = 1000\;\text{kg/m}^3$.
Danach: Lösung scannen oder fotografieren (ohne Namen und Standort) und mit dem \emph{Bewertungspaket} bewerten lassen.
\end{lpinfo}

\lpteil{Teil A · Druck}{Kapitel 1 · 4 Punkte}

\begin{lpaufgabe}{G1}{Raupenfahrzeug im Diagramm}{4 P}
Das Diagramm zeigt, welchen Druck ein Raupenfahrzeug auf den Boden ausübt, je nachdem, wie gross seine Auflagefläche \(A\) ist (verschieden breite Raupen).\par\smallskip
\begin{center}\begin{tikzpicture}[x=1.55cm,y=0.024cm,>=latex]
\foreach \x in {0,0.5,...,8} \draw[lplinie!70,very thin] (\x,0) -- (\x,160);
\foreach \y in {0,5,...,160} \draw[lplinie!70,very thin] (0,\y) -- (8,\y);
\foreach \y in {0,20,...,160} {\draw[lplinie] (0,\y) -- (8,\y); \node[left] at (0,\y) {\small \y};}
\foreach \x in {0,...,8} {\draw[lplinie] (\x,0) -- (\x,160); \node[below] at (\x,0) {\small \x};}
\draw[->] (0,0) -- (8.4,0) node[below left,yshift=-12pt] {\small \(A\) in \(\text{m}^2\)};
\draw[->] (0,0) -- (0,172) node[right] {\small \(p\) in kPa};
\draw[very thick,lpbernstein] plot[domain=0.94:8,samples=120] (\x,{150/\x});
\end{tikzpicture}\end{center}
\begin{enumerate}[label=\alph*),leftmargin=*,itemsep=2pt,topsep=2pt]
\item Lies ab: Welchen Druck übt das Fahrzeug mit \(3\;\text{m}^2\) Auflagefläche aus?
\item Berechne aus diesem Wertepaar die Gewichtskraft und die Masse des Fahrzeugs.
\item Ein Moorboden trägt höchstens \(25\;\text{kPa}\). Lies ab, wie gross die Auflagefläche mindestens sein muss.
\item Begründe, warum die Kurve nach rechts immer flacher wird und die Achse nie erreicht.
\end{enumerate}
\lpplatz{8}
\end{lpaufgabe}

\lpteil{Teil B · Schweredruck und Luftdruck}{Kapitel 2 und 3 · 8 Punkte}

\begin{lpaufgabe}{G2}{Tauchen im Bergsee}{4 P}
Eine Taucherin ist in einem Bergsee, an dessen Oberfläche der Luftdruck \(840\;\text{hPa}\) beträgt. Ihr Tauchcomputer zeigt einen Gesamtdruck von \(2.30\;\text{bar}\).
\begin{enumerate}[label=\alph*),leftmargin=*,itemsep=2pt,topsep=2pt]
\item Wie tief ist sie?
\item In welcher Tiefe zeigt derselbe Tauchcomputer \(2.30\;\text{bar}\) im Meer (\(1025\;\text{kg/m}^3\), Luftdruck \(1013\;\text{hPa}\))?
\item Begründe, warum derselbe Gesamtdruck im Bergsee in grösserer Tiefe erreicht wird.
\end{enumerate}
\lpplatz{7}
\end{lpaufgabe}

\begin{lpaufgabe}{G3}{Pascals Fassversuch}{4 P}
Blaise Pascal soll diesen Versuch gezeigt haben: Ein Holzfass ist \(1.6\;\text{m}\) hoch und randvoll mit Wasser. In den Deckel steckt er ein dünnes, \(8\;\text{m}\) langes Rohr und füllt es mit wenigen Litern Wasser. Das Fass platzt.
\begin{enumerate}[label=\alph*),leftmargin=*,itemsep=2pt,topsep=2pt]
\item Wie gross ist der Schweredruck am Fassboden, bevor das Rohr gefüllt ist?
\item Wie gross ist er, wenn das Rohr bis oben gefüllt ist?
\item Wenige Liter Wasser — und das Fass platzt. Begründe, warum die kleine Wassermenge im Rohr genügt.
\end{enumerate}
\lpplatz{7}
\end{lpaufgabe}

\lpteil{Teil C · Pascal und Archimedes}{Kapitel 4, 5 und 6 · 13 Punkte}

\begin{lpaufgabe}{G4}{Wagenheber}{4 P}
Ein hydraulischer Wagenheber soll ein Auto (\(1200\;\text{kg}\)) heben; sein grosser Kolben hat \(400\;\text{cm}^2\). Mit der Hand kannst du am Pumpkolben bequem \(120\;\text{N}\) aufbringen.
\begin{enumerate}[label=\alph*),leftmargin=*,itemsep=2pt,topsep=2pt]
\item Welcher Druck herrscht im Öl, in bar?
\item Wie gross darf die Fläche des Pumpkolbens höchstens sein, in \(\text{cm}^2\)?
\item Das Auto soll \(12\;\text{cm}\) gehoben werden. Wie viel Öl (in Liter) muss dafür hinübergepumpt werden, und welchen Weg legt der Pumpkolben insgesamt zurück (über alle Pumpstösse)?
\end{enumerate}
\lpplatz{7}
\end{lpaufgabe}

\begin{lpaufgabe}{G5}{Ballon}{5 P}
Ein Kinderballon fasst \(12\;\text{l}\) Helium; Hülle und Gas haben zusammen \(8\;\text{g}\). Das archimedische Prinzip gilt auch in Luft (\(\rho_\text{Luft} = 1.20\;\text{kg/m}^3\)).
\begin{enumerate}[label=\alph*),leftmargin=*,itemsep=2pt,topsep=2pt]
\item Skizziere den Ballon mit den Kräften, die auf ihn wirken, mit Namen und Richtung.
\item Wie gross ist der Auftrieb in Luft?
\item Wie gross ist die Gewichtskraft von Hülle und Gas?
\item Welche Masse kann der Ballon zusätzlich tragen, ohne zu sinken?
\item Derselbe Ballon wird unter Wasser festgebunden. Ist der Auftrieb dort grösser, gleich gross oder kleiner als in Luft? Begründe.
\end{enumerate}
\skizzenfeld{2.6cm}
\lpplatz{5}
\end{lpaufgabe}

\begin{lpaufgabe}{G6}{Aräometer}{4 P}
Ein Aräometer ist ein Dichtemesser: ein beschwertes Glasröhrchen, das aufrecht in der Flüssigkeit schwimmt. Je nach Dichte taucht es verschieden tief ein. Betrachte es als schlanken Zylinder mit der Masse \(30\;\text{g}\) und dem Querschnitt \(2\;\text{cm}^2\).
\begin{enumerate}[label=\alph*),leftmargin=*,itemsep=2pt,topsep=2pt]
\item Wie tief taucht es in Wasser ein?
\item Wie tief taucht es in Spiritus (\(790\;\text{kg/m}^3\)) ein?
\item Auf der Skala des Aräometers stehen die grossen Dichten unten, die kleinen oben. Begründe.
\item In einer unbekannten Flüssigkeit taucht es \(13.6\;\text{cm}\) tief ein. Welche Dichte hat sie?
\end{enumerate}
\lpplatz{7}
\end{lpaufgabe}

\vfill
{\small\color{lpgrau}Bewerten: mit dem Bewertungspaket (Musterlösung, Punkteraster, Auftrag an eine KI) — erst nach dem Lösen öffnen.}
\end{document}
